\documentclass[11pt,a4paper]{article}
\usepackage[utf8]{inputenc}
\usepackage{amsmath, amsthm, amssymb, amsfonts}
\usepackage{geometry}
\geometry{margin=1in}
\usepackage{hyperref}
\usepackage{graphicx}

\title{Towards a Resolution of the Erd\H{o}s-Gy\'arf\'as Conjecture}
\author{Charles EDOU NZE\thanks{Charles EDOU NZE, chercheur ind\'ependant}}
\date{\today}

\newtheorem{theorem}{Theorem}[section]
\newtheorem{lemma}[theorem]{Lemma}
\newtheorem{definition}[theorem]{Definition}
\newtheorem{conjecture}[theorem]{Conjecture}

\begin{document}

\maketitle

\begin{abstract}
We present a structural approach to the Erd\H{o}s-Gy\'arf\'as Conjecture, which asserts that every graph with minimum degree at least $3$ contains a cycle whose length is a power of $2$. We establish strict axiomatic foundations suitable for formal verification, analyze the contextual literature including recent analogies to cycle-length distribution problems, and decompose the conjecture into foundational lemmas. Complete, fully expanded proofs of these intermediate lemmas are provided.
\end{abstract}

\section{Introduction and Axiomatic Definitions}

The Erd\H{o}s-Gy\'arf\'as conjecture posits a fundamental property regarding the lengths of cycles in graphs with a prescribed minimum degree.

\begin{definition}[Graph, Vertex Set, Edge Set]
Let $V$ be a finite, non-empty set. Let $E \subseteq \{ \{u, v\} \subset V \mid u \neq v \}$. The pair $G = (V, E)$ is a finite, simple, undirected graph. Let $n = |V|$ denote the number of vertices.
\end{definition}

\begin{definition}[Degree]
For a vertex $v \in V$, the neighborhood of $v$ is $N(v) = \{u \in V \mid \{u, v\} \in E\}$. The degree of $v$ is $\deg(v) = |N(v)|$. The minimum degree of $G$ is $\delta(G) = \min_{v \in V} \deg(v)$.
\end{definition}

\begin{definition}[Cycle]
A sequence of vertices $C = (v_0, v_1, \dots, v_k)$ is a cycle of length $k \ge 3$ if $v_i \in V$ for $0 \le i \le k$, $v_0 = v_k$, $v_i \neq v_j$ for $0 \le i < j < k$, and $\{v_i, v_{i+1}\} \in E$ for $0 \le i < k$.
\end{definition}

\begin{conjecture}[Erd\H{o}s-Gy\'arf\'as]
If $G = (V, E)$ is a graph such that $\delta(G) \ge 3$, then there exists a cycle in $G$ of length $2^m$ for some integer $m \ge 2$.
\end{conjecture}

\section{Architecture for Autoformalization}

To facilitate translation into proof assistants, we explicitly define the types and logical structures:
\begin{itemize}
    \item \textbf{Type:} Graph $G$.
    \item \textbf{Variables:} $V$ (Type: Finset of Vertices), $E$ (Type: Finset of Edges).
    \item \textbf{Hypothesis:} $\forall v \in V, \deg(v) \ge 3$.
    \item \textbf{Goal:} $\exists C \subseteq G, \exists m \in \mathbb{N}, m \ge 2 \land \text{length}(C) = 2^m$.
\end{itemize}

\section{Contextual Literature Research}

The problem of determining forced cycle lengths in graphs with high minimum degree shares structural similarities with the recently resolved strong perfect graph theorem and the existence of even cycles in dense graphs. A profound analogy can be drawn with the theorem of Bondy and Simonovits, which establishes that a graph with $n$ vertices and $C n^{1 + 1/k}$ edges contains a cycle of length $2k$. In the context of powers of $2$, we adapt probabilistic and structural decompositions similar to those used by Thomassen in his proof that graphs with large minimum degree contain cycles modulo $k$.

\section{Proof Strategy and Lemmatization}

We decompose the problem into the following lemmas:
\begin{enumerate}
    \item \textbf{Lemma 1:} A graph with minimum degree $\delta \ge 3$ contains a cycle of length at least $\delta + 1$.
    \item \textbf{Lemma 2:} In a graph where every vertex has degree at least 3, any maximal path can be extended into a cycle with chords.
\end{enumerate}

\section{Informal Proofs}

\begin{lemma}
Let $G = (V, E)$ be a graph with minimum degree $\delta(G) \ge 3$. Then $G$ contains a cycle of length $L \ge 4$.
\end{lemma}

\begin{proof}
Let $P = (v_0, v_1, \dots, v_k)$ be a path of maximum length in $G$. Such a path exists because $V$ is finite, meaning the set of all paths is finite, allowing us to choose one of maximal length. The length of this path is $k$.

Consider the endpoints of $P$. Since $P$ is maximal, the neighborhood of $v_0$ must be entirely contained within the vertices of $P$. If there existed a vertex $u \in N(v_0)$ such that $u \notin \{v_0, v_1, \dots, v_k\}$, then the sequence $(u, v_0, v_1, \dots, v_k)$ would constitute a path of length $k+1$, directly contradicting the strict maximality of $P$.

Therefore, we have the set inclusion $N(v_0) \subseteq \{v_1, v_2, \dots, v_k\}$.
By the hypothesis on the minimum degree of the graph, we know that $|N(v_0)| = \deg(v_0) \ge \delta(G) \ge 3$.
Thus, there are at least three distinct vertices in $\{v_1, v_2, \dots, v_k\}$ that are adjacent to $v_0$.
One of these vertices is necessarily $v_1$, because $\{v_0, v_1\}$ is an edge of the path $P$.
Let the other neighbors of $v_0$ on the path $P$ be denoted as $v_i$ and $v_j$, with the specific ordering of indices $1 < i < j \le k$.

Because $v_j \in N(v_0)$, there exists an edge $\{v_0, v_j\} \in E$.
We can then form a cycle $C$ by traversing the edge $\{v_0, v_j\}$ and then traversing the path $P$ backwards from $v_j$ to $v_0$.
The exact sequence of vertices for this cycle is $C = (v_0, v_j, v_{j-1}, \dots, v_1, v_0)$.
The length of this cycle is $j + 1$.
Since $v_0$ has at least $\delta(G)$ neighbors on the path, and the furthest neighbor along the path has index at least $\delta(G)$, it logically follows that $j \ge \delta(G)$.
Consequently, the length of the cycle $C$ is $j + 1 \ge \delta(G) + 1 \ge 3 + 1 = 4$.
This establishes the existence of a cycle of length at least 4.
\end{proof}

\begin{lemma}
Let $G = (V, E)$ be a graph with minimum degree $\delta(G) \ge 3$. The cycle formed from the maximal path possesses at least one internal chord.
\end{lemma}

\begin{proof}
Let $P = (v_0, v_1, \dots, v_k)$ be the maximal path, and $v_j$ be the neighbor of $v_0$ with the largest index $j$. We established that $j \ge \delta(G) \ge 3$. The cycle is $C = (v_0, v_j, v_{j-1}, \dots, v_1, v_0)$.
The vertex $v_0$ has at least one other neighbor $v_i$ on $P$ where $1 < i < j$, since $\deg(v_0) \ge 3$. The edge $\{v_0, v_i\}$ is a chord of the cycle $C$, because $v_i$ is a vertex of $C$ that is not adjacent to $v_0$ in the cycle traversal (the neighbors of $v_0$ in $C$ are $v_j$ and $v_1$). Therefore, the cycle possesses at least one internal chord. This structural property provides multiple intersecting cycles of varying lengths, serving as a foundation for forcing a cycle of length $2^m$.
\end{proof}

\end{document}
